% !TeX program = lualatex
% ================================================================
% Polish Mathematical Dictionary
%
% Generated by the server using the following settings:
%
%   language            Polish (pol)
%   text direction      left to right
%   font                Noto Serif
%   fallback font       Noto Serif
%   built from          latex/dictionary/mathematicalDictionary.tex
%   tier 3 vocabulary   green   RGB 0 128 0
%   English text        black   RGB 0 0 0
%   Polish text         purple  RGB 112 48 160
%   layout macros       version 2
%
% Please don't edit any information in this comment block - the server
% compares it against the current settings and rebuilds this file when
% they differ, so changes here would be overwritten. However, feel free
% to make updates to the LaTeX below if you want to edit it.
% ================================================================
%
% ---- What is safe to edit in this file ----------------------------------
%
% Safe:     the Polish word and the Polish definition - the third
%           and fourth arguments of each entry. That is what this file is for:
%           if a translation reads wrongly to somebody who speaks the language,
%           correct it here and every sheet that uses the word follows.
%
% Not safe: the English word and the English definition - the first and second
%           arguments. The first is how an entry is matched to the shared
%           dictionary; the second is the wording the translation was made
%           from, and is how the server knows the translation is still current.
%           Changing an English definition here does not change the shared
%           dictionary - it only makes this entry look up to date when it is
%           not. Edit the English in the shared dictionary instead, and this
%           file will be brought back into step on the next run.
%
% Entries are added and removed by the server, following the shared dictionary.
% An entry the server cannot read is reported rather than guessed at, so if a
% run complains about this file, look for a missing brace.
% -------------------------------------------------------------------------

\documentclass[11pt]{article}
\usepackage[a4paper,margin=18mm]{geometry}
% Characters this language's font has not got are borrowed from another,
% rather than being silently dropped - see docs/translated-sheet-latex.md
\directlua{%
  if luaotfload and luaotfload.add_fallback then
    luaotfload.add_fallback("eallatinfallback",
      { "Noto Serif:mode=harf;" })
  else
    tex.error("this luaotfload is too old to register a fallback font")
  end
}
\usepackage[english]{babel}
\babelprovide[import]{polish}
\babelfont[polish]{rm}[Renderer=Harfbuzz, BoldFont={Noto Serif}, ItalicFont={Noto Serif}, BoldItalicFont={Noto Serif}, RawFeature={fallback=eallatinfallback}]{Noto Serif}
\babelfont[polish]{sf}[Renderer=Harfbuzz, BoldFont={Noto Serif}, ItalicFont={Noto Serif}, BoldItalicFont={Noto Serif}, RawFeature={fallback=eallatinfallback}]{Noto Serif}
\newcommand{\ealtext}[1]{\foreignlanguage{polish}{#1}}
\newcommand{\ealtextblock}[1]{{\raggedright\foreignlanguage{polish}{#1}\par}}
% ================================================================
% EAL layout helpers
% ================================================================
\usepackage{xcolor}

\definecolor{ealkeycolour}{RGB}{0,128,0}
\definecolor{ealenglish}{RGB}{0,0,0}
\definecolor{ealtwocolour}{RGB}{112,48,160}

\newsavebox{\ealboxen}
\newsavebox{\ealboxtr}
\newlength{\ealglosswd}
\newlength{\ealglossraise}
\setlength{\ealglossraise}{1.15em}

% a tier 3 word, in English and in translation alike
\newcommand{\ealkey}[1]{\textcolor{ealkeycolour}{#1}}
\newcommand{\ealkeytr}[1]{\textcolor{ealkeycolour}{\ealtext{#1}}}

% One translated word above one English word. Built from kernel
% primitives, never a tabular, which would break wherever the sheet
% loads array, booktabs or tabularx. The box is as wide as the wider of
% the two, so a long translation pushes its neighbours aside instead of
% overlapping them, and it claims its full height so a table row or a
% TikZ node makes room for it.
\newcommand{\ealgloss}[2]{%
  \sbox{\ealboxen}{\ealkey{#1}}%
  \sbox{\ealboxtr}{\small\textcolor{ealtwocolour}{\ealtext{#2}}}%
  \setlength{\ealglosswd}{\wd\ealboxen}%
  \ifdim\wd\ealboxtr>\ealglosswd\setlength{\ealglosswd}{\wd\ealboxtr}\fi
  \leavevmode
  \makebox[\ealglosswd][c]{%
    \makebox[0pt][c]{\usebox{\ealboxen}}%
    \makebox[0pt][c]{\raisebox{\ealglossraise}{\usebox{\ealboxtr}}}%
  }%
}

% opens up line spacing around a block of glosses, so the raised
% translations sit clear of the line above
\newenvironment{ealglossed}{\par\linespread{2.1}\selectfont\ignorespaces}{\par}

% a whole translated sentence above its English counterpart. block level,
% so it does not belong inside a tabular cell
\newcommand{\ealpara}[2]{%
  \par\addvspace{0.5\baselineskip}%
  {\color{ealtwocolour}\ealtextblock{#1}}%
  \penalty200
  {\color{ealenglish}#2\par}%
  \addvspace{0.5\baselineskip}%
}

% ================================================================
% Dictionary layout
% ================================================================
\usepackage{multicol}

\newlength{\ealdictgap}
\setlength{\ealdictgap}{1.2mm}

% english word / english meaning / translated word / translated meaning
\newcommand{\dictentrytr}[4]{%
  \par\addvspace{\ealdictgap}%
  \noindent
  {\bfseries\ealkey{#1}}\quad{\small #2}\par
  \nopagebreak
  {\color{ealtwocolour}\ealtext{{\bfseries #3}~\textemdash{} {\small #4}}\par}%
  \par
}

\pagestyle{empty}
\setlength{\parindent}{0pt}

\begin{document}
{\large\bfseries Polish Mathematical Dictionary}\par\medskip

\dictentrytr{absolute value}{the distance of a number from zero, ignoring whether it is positive or negative.}{wartość bezwzględna}{odległość liczby od zera, bez uwzględniania tego, czy jest ona dodatnia, czy ujemna.}
\dictentrytr{addition}{combining two or more numbers to find their total.}{dodawanie}{łączenie dwóch lub więcej liczb w celu znalezienia ich sumy.}
\dictentrytr{adjacent}{next to, and sharing a side or a corner}{sąsiedni}{znajdujący się obok i mający wspólny bok lub wierzchołek}
\dictentrytr{altitude}{height above sea level or another fixed reference level.}{wysokość}{wysokość nad poziomem morza lub innym stałym poziomem odniesienia.}
\dictentrytr{angle}{the amount of turn between two straight lines that meet at a point}{kąt}{wielkość obrotu między dwiema prostymi, które przecinają się w punkcie.}
\dictentrytr{annular}{Describing a shape made by the space between two concentric circles.}{pierścieniowy}{Opisujący kształt utworzony przez przestrzeń między dwoma współśrodkowymi okręgami.}
\dictentrytr{apex angle}{the angle at the point of a triangle opposite its base.}{kąt przy wierzchołku}{kąt przy wierzchołku trójkąta leżącym naprzeciwko jego podstawy.}
\dictentrytr{approximately}{close to but not exactly; about.}{w przybliżeniu}{blisko, ale nie dokładnie; około.}
\dictentrytr{arc}{part of the curved edge of a circle}{łuk}{część zakrzywionej krawędzi okręgu}
\dictentrytr{arccos}{The operation that finds an angle when its cosine is known.}{arcus cosinus}{Operacja, która znajduje kąt, gdy znany jest jego cosinus.}
\dictentrytr{arctan}{The operation that finds an angle when its tangent is known.}{arcus tangens}{Operacja, która znajduje kąt, gdy znany jest jego tangens.}
\dictentrytr{area}{the amount of space inside a flat shape}{pole}{ilość miejsca wewnątrz figury płaskiej}
\dictentrytr{area model}{a rectangle diagram split into smaller rectangles, used to show how to multiply fractions.}{model pola}{diagram w kształcie prostokąta podzielony na mniejsze prostokąty, używany do pokazania, jak mnożyć ułamki.}
\dictentrytr{arithmetic}{the branch of mathematics that deals with adding, subtracting, multiplying and dividing numbers.}{arytmetyka}{dział matematyki zajmujący się dodawaniem, odejmowaniem, mnożeniem i dzieleniem liczb.}
\dictentrytr{ascending}{arranged from smallest to largest.}{rosnąco}{ułożone od najmniejszego do największego.}
\dictentrytr{ascending order}{arranged from smallest to largest}{kolejność rosnąca}{uporządkowana od najmniejszego do największego.}
\dictentrytr{average}{a single value chosen to stand for a whole set of data}{średnia}{jedna wartość wybrana do reprezentowania całego zestawu danych}
\dictentrytr{average speed}{total distance travelled divided by total time taken.}{prędkość średnia}{całkowita przebyta droga podzielona przez całkowity czas.}
\dictentrytr{axis}{one of the lines a graph is measured against}{oś}{jedna z linii, w odniesieniu do których mierzony jest wykres}
\dictentrytr{bar model}{a diagram made of rectangles that shows the parts and total of an equation.}{model słupkowy}{diagram złożony z prostokątów, który pokazuje części i całość równania.}
\dictentrytr{base}{the side of a shape chosen as the bottom for working out its area.}{podstawa}{bok figury wybrany jako podstawa do obliczania jej pola.}
\dictentrytr{basis point}{a unit equal to one hundredth of one percent, used for small changes in rates.}{punkt bazowy}{jednostka równa jednej setnej procenta, używana do małych zmian stóp procentowych.}
\dictentrytr{bearing}{a direction given as an angle clockwise from north, written with three figures}{azymut}{kierunek podany jako kąt mierzony zgodnie z ruchem wskazówek zegara od północy, zapisywany za pomocą trzech cyfr}
\dictentrytr{bimodal}{having two values that appear most often in a list of numbers.}{dwumodalny}{mający dwie wartości, które występują najczęściej na liście liczb.}
\dictentrytr{bisector}{a line that cuts an angle or another line exactly in half}{dwusieczna}{linia, która dzieli kąt lub inną linię dokładnie na pół}
\dictentrytr{bodmas}{a rule that says which operation to do first in a calculation.}{kolejność wykonywania działań}{reguła, która mówi, które działanie wykonać najpierw w obliczeniach.}
\dictentrytr{bond}{a type of loan to a government or company that pays regular money and is paid back later.}{obligacja}{rodzaj pożyczki udzielonej rządowi lub firmie, która regularnie wypłaca pieniądze i jest później spłacana.}
\dictentrytr{box method}{a way of multiplying or expanding by splitting a rectangle into smaller parts.}{metoda pudełkowa}{sposób mnożenia lub rozwijania wyrażeń przez podzielenie prostokąta na mniejsze części.}
\dictentrytr{bracket}{one of a pair of curved symbols, ( ), used to group numbers or letters.}{nawias}{jeden z pary zakrzywionych symboli ( ), używany do grupowania liczb lub liter.}
\dictentrytr{brackets}{A pair of signs used in mathematics to group together a part of an expression.}{nawiasy}{para znaków używana w matematyce do grupowania części wyrażenia.}
\dictentrytr{bus stop method}{a written method for division using a roof-like sign and a divisor.}{dzielenie pisemne}{pisemna metoda dzielenia z użyciem znaku w kształcie daszka i dzielnika.}
\dictentrytr{cancel}{to remove a shared number from the top and bottom of a fraction before multiplying.}{skracać}{usuwać wspólną liczbę z licznika i mianownika ułamka przed mnożeniem.}
\dictentrytr{cardinal directions}{The four main compass points: North, South, East and West.}{strony świata}{cztery główne kierunki kompasu: północ, południe, wschód i zachód.}
\dictentrytr{category}{a group of non-numerical items, for example a colour, used as data.}{kategoria}{grupa elementów nienumerycznych, na przykład kolor, używana jako dane.}
\dictentrytr{centre}{the fixed middle point of a circle}{środek}{stały punkt środkowy koła.}
\dictentrytr{chord}{a straight line joining two points on the edge of a circle}{cięciwa}{odcinek łączący dwa punkty na okręgu}
\dictentrytr{chunking}{a division method that breaks a larger number into easier smaller parts.}{metoda chunking}{metoda dzielenia, która rozkłada większą liczbę na prostsze, mniejsze części.}
\dictentrytr{circle}{a round flat shape whose edge is always the same distance from its centre}{koło}{okrągły płaski kształt, którego brzeg jest zawsze w tej samej odległości od środka.}
\dictentrytr{circular}{round in shape, with every point on its edge the same distance from the centre.}{okrągły}{mający kształt koła; każdy punkt na jego brzegu jest w tej samej odległości od środka.}
\dictentrytr{circumference}{the distance all the way round a circle}{obwód}{odległość dookoła całego okręgu}
\dictentrytr{coefficient}{the number in front of a letter, showing how many of it there are}{współczynnik}{liczba stojąca przed literą, pokazująca, ile jest tych liter}
\dictentrytr{column vector}{a vector written vertically with one number above another.}{wektor kolumnowy}{wektor zapisany pionowo, z jedną liczbą nad drugą.}
\dictentrytr{common denominator}{a bottom number shared by two or more fractions, used for adding or comparing them.}{wspólny mianownik}{liczba w mianowniku wspólna dla dwóch lub więcej ułamków, używana do ich dodawania lub porównywania.}
\dictentrytr{common difference}{the fixed number added to each term to get the next term}{różnica ciągu arytmetycznego}{stała liczba dodawana do każdego wyrazu, aby otrzymać następny wyraz.}
\dictentrytr{common factor}{a number that divides exactly into two or more other numbers.}{wspólny dzielnik}{liczba, która dzieli bez reszty dwie lub więcej innych liczb.}
\dictentrytr{common ratio}{the fixed number used to create each next term in a geometric sequence}{iloraz ciągu geometrycznego}{stała liczba używana do utworzenia każdego następnego wyrazu w ciągu geometrycznym.}
\dictentrytr{commutativity}{the rule that swapping two numbers before you add or times them leaves the answer the same}{przemienność}{zasada, że zamiana miejscami dwóch liczb przed dodaniem ich lub pomnożeniem pozostawia wynik taki sam.}
\dictentrytr{component}{one of the numbers that make up a vector.}{składowa}{jedna z liczb tworzących wektor.}
\dictentrytr{compounding}{the process of earning profit on both the original amount and profit already earned.}{kapitalizacja}{proces uzyskiwania zysku zarówno od pierwotnej kwoty, jak i od zysku już osiągniętego.}
\dictentrytr{concentric}{Having the same centre.}{współśrodkowy}{mający ten sam środek.}
\dictentrytr{congruent}{exactly the same shape and size}{przystający}{dokładnie ten sam kształt i ta sama wielkość}
\dictentrytr{constant}{a fixed number that stands alone, like 3 in 3x+3.}{stała}{stała liczba, która stoi samodzielnie, na przykład 3 w wyrażeniu 3x+3.}
\dictentrytr{convergence}{the process of getting closer and closer to a particular value}{zbieżność}{proces zbliżania się coraz bardziej do określonej wartości.}
\dictentrytr{conversion}{a change of a value from one form, such as a fraction, to another, such as a percentage.}{zamiana}{zmiana wartości z jednej postaci, na przykład ułamka, na inną, na przykład procent.}
\dictentrytr{convert}{to change a number from one form into another, such as a mixed number into an improper fraction}{zamieniać}{zmieniać liczbę z jednej postaci na inną, np. liczbę mieszaną na ułamek niewłaściwy.}
\dictentrytr{coordinate}{a number or pair of numbers that fixes the position of a point.}{współrzędna}{liczba lub para liczb, która ustala położenie punktu.}
\dictentrytr{correlation}{how closely two sets of data are related}{korelacja}{stopień, w jakim dwa zestawy danych są ze sobą powiązane}
\dictentrytr{cos}{the adjacent side divided by the hypotenuse for an angle in a right triangle.}{cosinus}{przyprostokątna przyległa do kąta podzielona przez przeciwprostokątną dla kąta w trójkącie prostokątnym.}
\dictentrytr{cosine}{in a three-sided shape with a right angle, the ratio of the adjacent side to the hypotenuse.}{cosinus}{w trójkącie prostokątnym stosunek przyprostokątnej przyległej do przeciwprostokątnej.}
\dictentrytr{cosine rule}{a formula connecting three side lengths and one angle in a three-sided shape.}{twierdzenie cosinusów}{wzór łączący długości trzech boków i jeden kąt w trójkącie.}
\dictentrytr{counterexample}{an example that proves a general statement is not always true}{kontrprzykład}{przykład, który dowodzi, że ogólne stwierdzenie nie zawsze jest prawdziwe.}
\dictentrytr{coupon}{The regular interest payment made on a bond.}{kupon}{regularna płatność odsetkowa dokonywana z tytułu obligacji.}
\dictentrytr{cube}{the number you get when a value is multiplied by itself twice.}{sześcian}{liczba, którą otrzymujesz, gdy wartość jest mnożona przez siebie dwukrotnie.}
\dictentrytr{cuboid}{a 3D solid shape whose six faces are all rectangles.}{prostopadłościan}{bryła przestrzenna, której sześć ścian jest prostokątami.}
\dictentrytr{cumulative}{adding up as you go along, so each total includes the ones before}{skumulowany}{dodający się w miarę postępu, tak że każda suma obejmuje poprzednie}
\dictentrytr{data}{the set of values collected}{dane}{zbiór zebranych wartości}
\dictentrytr{dataset}{a collection of numbers or information being studied}{zbiór danych}{zbiór liczb lub informacji poddawanych badaniu.}
\dictentrytr{decimal}{a number written with a point, where the digits after the point are parts of a whole}{ułamek dziesiętny}{liczba zapisana z przecinkiem, gdzie cyfry po przecinku są częściami całości}
\dictentrytr{decimal place}{the position of a digit after the decimal point, such as tenths or hundredths}{miejsce dziesiętne}{położenie cyfry po przecinku, na przykład części dziesiąte lub setne.}
\dictentrytr{decimal places}{the digits written after the point in a number.}{miejsca dziesiętne}{cyfry zapisane po przecinku w liczbie.}
\dictentrytr{decimal point}{the dot used to separate the whole part of a number from the fractional part.}{przecinek dziesiętny}{przecinek używany do oddzielania części całkowitej liczby od części ułamkowej.}
\dictentrytr{degree}{a unit for measuring angles; 360 of them make one full turn.}{stopień}{jednostka miary kątów; 360 z nich tworzy jeden pełny obrót.}
\dictentrytr{denominator}{the number below the line in a fraction, showing how many equal parts the whole is split into}{mianownik}{liczba pod kreską ułamkową, pokazująca, na ile równych części podzielono całość}
\dictentrytr{depreciation}{the loss in value of something over time.}{amortyzacja}{spadek wartości czegoś w czasie.}
\dictentrytr{diameter}{a straight line across a circle through its centre}{średnica}{odcinek łączący dwa punkty na okręgu i przechodzący przez jego środek}
\dictentrytr{difference}{the answer when one number is subtracted from another}{różnica}{wynik odejmowania jednej liczby od drugiej}
\dictentrytr{digit}{a single figure from 0 to 9 that a number is written with}{cyfra}{pojedynczy znak od 0 do 9, za pomocą którego zapisuje się liczbę}
\dictentrytr{dimension}{a measurement of a shape, such as its length or width.}{wymiar}{miara figury, taka jak jej długość lub szerokość.}
\dictentrytr{dimensions}{the measurements of a shape, such as length, width and height.}{wymiary}{miary figury, takie jak długość, szerokość i wysokość.}
\dictentrytr{direct proportion}{a relationship where two amounts change together while their ratio stays the same.}{proporcjonalność prosta}{zależność, w której dwie wielkości zmieniają się razem, a ich stosunek pozostaje taki sam.}
\dictentrytr{directed path}{a line with an arrow showing direction from one place to another}{ścieżka skierowana}{linia ze strzałką pokazująca kierunek z jednego miejsca do drugiego.}
\dictentrytr{directly proportional}{when one amount is multiplied by a number, the other amount is multiplied by that same number.}{wprost proporcjonalny}{gdy jedna wielkość jest mnożona przez liczbę, druga wielkość jest mnożona przez tę samą liczbę.}
\dictentrytr{discriminant}{a value calculated from the numbers in an equation; if it is negative, the equation has no real solutions.}{wyróżnik}{wartość obliczana z liczb w równaniu; jeśli jest ujemna, równanie nie ma rozwiązań rzeczywistych.}
\dictentrytr{displacement}{the straight-line distance and direction from a starting point to a final position.}{przemieszczenie}{odległość w linii prostej i kierunek od punktu początkowego do położenia końcowego.}
\dictentrytr{distributive property}{the rule that a(b+c)=ab+ac.}{rozdzielność mnożenia względem dodawania}{zasada, że a(b+c)=ab+ac.}
\dictentrytr{divide}{to split a number into equal parts, or find how many times one number fits inside another.}{dzielić}{podzielić liczbę na równe części lub ustalić, ile razy jedna liczba mieści się w drugiej.}
\dictentrytr{division}{the operation of splitting a number into equal parts.}{dzielenie}{działanie polegające na podzieleniu liczby na równe części.}
\dictentrytr{divisor}{the number that another number is divided by.}{dzielnik}{liczba, przez którą dzielona jest inna liczba.}
\dictentrytr{duration}{the length of time until a bond or gilt is paid back.}{czas trwania}{długość czasu do spłaty obligacji lub obligacji skarbowej.}
\dictentrytr{edge}{a straight line where two faces of a solid meet}{krawędź}{linia prosta, w której spotykają się dwie ściany bryły}
\dictentrytr{elimination}{removing one unknown from simultaneous equations by adding or subtracting them}{metoda eliminacji}{usuwanie jednej niewiadomej z układu równań przez dodawanie lub odejmowanie tych równań.}
\dictentrytr{embedding}{a way of turning words into coordinates or vectors so meanings can be compared mathematically}{osadzenie}{sposób zamiany słów na współrzędne lub wektory, dzięki czemu znaczenia można porównywać matematycznie.}
\dictentrytr{endpoint}{the point at which a route or line finishes}{punkt końcowy}{punkt, w którym trasa lub linia się kończy.}
\dictentrytr{equal}{having the same size or number}{równy}{mający ten sam rozmiar lub liczbę.}
\dictentrytr{equation}{a mathematical statement with an equals sign showing that two sides are equal.}{równanie}{matematyczne stwierdzenie ze znakiem równości pokazujące, że obie strony są równe.}
\dictentrytr{equilateral}{having all sides the same length}{równoboczny}{mający wszystkie boki tej samej długości.}
\dictentrytr{equilateral triangle}{a triangle with all three sides the same length.}{trójkąt równoboczny}{trójkąt, którego wszystkie trzy boki mają tę samą długość.}
\dictentrytr{equivalent}{having the same value, even though it is written differently}{równoważny}{mający tę samą wartość, nawet jeśli zapisany inaczej}
\dictentrytr{equivalent fraction}{a different fraction that means the same amount as another fraction.}{ułamek równoważny}{inny ułamek, który oznacza tę samą ilość, co inny ułamek.}
\dictentrytr{estimate}{a sensible answer worked out quickly from rounded numbers}{oszacowanie}{rozsądna odpowiedź szybko obliczona z zaokrąglonych liczb}
\dictentrytr{evaluate}{to work out the value of an expression}{obliczyć}{obliczyć wartość wyrażenia.}
\dictentrytr{even}{able to be divided exactly by 2.}{parzysty}{dający się dokładnie podzielić przez 2.}
\dictentrytr{even number}{a whole number that can be divided exactly by two.}{liczba parzysta}{liczba całkowita, którą można podzielić dokładnie przez dwa.}
\dictentrytr{expand}{to multiply out brackets so the expression has no brackets left}{wymnażać nawiasy}{wymnożyć nawiasy, tak aby wyrażenie nie miało już nawiasów}
\dictentrytr{exponent}{the raised number showing how many times to multiply a number by itself.}{wykładnik}{liczba w indeksie górnym pokazująca, ile razy pomnożyć liczbę przez siebie.}
\dictentrytr{expression}{a group of numbers and letters joined by operations, with no equals sign}{wyrażenie}{grupa liczb i liter połączonych działaniami, bez znaku równości}
\dictentrytr{face}{a flat surface of a solid shape}{ściana}{płaska powierzchnia bryły}
\dictentrytr{factor}{a whole number that divides exactly into another number}{dzielnik}{liczba całkowita, która dzieli bez reszty inną liczbę}
\dictentrytr{factorise}{to write an expression as a product, usually by putting brackets back in}{rozkładać na czynniki}{zapisać wyrażenie jako iloczyn, zwykle przez wstawienie nawiasów z powrotem}
\dictentrytr{fibonacci}{a famous sequence where each number is the total of the two before it}{ciąg Fibonacciego}{słynny ciąg, w którym każda liczba jest sumą dwóch poprzednich.}
\dictentrytr{finite}{limited in size; not going on forever.}{skończony}{ograniczony co do wielkości; nie trwający w nieskończoność.}
\dictentrytr{fixed point}{a value that stays the same when an iterative process is applied to it.}{punkt stały}{wartość, która pozostaje taka sama, gdy zastosuje się do niej proces iteracyjny.}
\dictentrytr{foil}{a way to multiply two brackets that each have two terms, multiplying First, Outer, Inner and Last.}{metoda FOIL}{sposób mnożenia dwóch nawiasów, z których każdy ma dwa wyrazy: mnoży się pierwsze, zewnętrzne, wewnętrzne i ostatnie wyrazy.}
\dictentrytr{formula}{a rule written with letters, showing how to work one quantity out from others}{wzór}{reguła zapisana literami, pokazująca, jak obliczyć jedną wielkość na podstawie innych}
\dictentrytr{fraction}{a number written as one whole split into equal parts, with a number of those parts taken}{ułamek}{liczba zapisana jako całość podzielona na równe części, z których wzięto pewną liczbę}
\dictentrytr{fraction wall}{a diagram of strips divided into equal parts, used to compare fractions.}{ściana ułamkowa}{diagram pasków podzielonych na równe części, używany do porównywania ułamków.}
\dictentrytr{frequency}{how many times something happens}{częstość}{ile razy coś się zdarza}
\dictentrytr{geometric}{describing a sequence where each term is multiplied by the same number to get the next.}{geometryczny}{opisujący ciąg, w którym każdy wyraz jest mnożony przez tę samą liczbę, aby otrzymać następny.}
\dictentrytr{geometry}{the branch of mathematics that studies shapes, sizes, and space.}{geometria}{dział matematyki, który bada kształty, rozmiary i przestrzeń.}
\dictentrytr{gilt}{A bond issued by the UK government to borrow money.}{brytyjska obligacja skarbowa}{obligacja wyemitowana przez rząd Wielkiej Brytanii w celu pożyczenia pieniędzy.}
\dictentrytr{gradient}{how steep a line is, found by how far it rises for each step across}{współczynnik kierunkowy}{jak stroma jest prosta, obliczany przez to, jak bardzo wznosi się na każdy krok w bok}
\dictentrytr{height}{the distance straight up from the base of a shape to its top, used to find area.}{wysokość}{odległość prosto w górę od podstawy figury do jej najwyższego punktu, używana do obliczania pola.}
\dictentrytr{hexagon}{a flat shape with six straight sides}{sześciokąt}{figura płaska o sześciu prostych bokach.}
\dictentrytr{highest common factor}{the largest whole number that divides exactly into two or more numbers}{największy wspólny dzielnik}{największa liczba całkowita, która dzieli bez reszty dwie lub więcej liczb}
\dictentrytr{hundredths}{the value of one part when a whole is divided into one hundred equal parts.}{setne}{wartość jednej części, gdy całość jest podzielona na sto równych części.}
\dictentrytr{hypotenuse}{the longest side of a right-angled triangle, opposite the right angle}{przeciwprostokątna}{najdłuższy bok trójkąta prostokątnego, leżący naprzeciw kąta prostego}
\dictentrytr{identity}{a statement that is true for every value of the letters in it}{tożsamość}{stwierdzenie, które jest prawdziwe dla każdej wartości liter w nim występujących}
\dictentrytr{improper fraction}{a fraction with a top number larger than or equal to its bottom number}{ułamek niewłaściwy}{ułamek, w którym licznik jest większy lub równy mianownikowi.}
\dictentrytr{index}{the small raised number showing what power something is raised to}{wykładnik}{mała liczba u góry pokazująca, do jakiej potęgi coś jest podniesione}
\dictentrytr{inequality}{a statement that one quantity is larger or smaller than another}{nierówność}{stwierdzenie, że jedna wielkość jest większa lub mniejsza od innej}
\dictentrytr{integer}{a whole number, which may be positive, negative or zero}{liczba całkowita}{liczba całkowita, która może być dodatnia, ujemna lub równa zero}
\dictentrytr{intercept}{the point where a graph crosses an axis}{punkt przecięcia z osią}{punkt, w którym wykres przecina oś}
\dictentrytr{interest}{Money paid for borrowing or earned on savings.}{odsetki}{pieniądze płacone za pożyczanie lub zarabiane na oszczędnościach.}
\dictentrytr{interest rate}{the percentage of an amount that is paid or earned as interest.}{stopa procentowa}{procent kwoty, który jest płacony lub otrzymywany jako odsetki.}
\dictentrytr{intersect}{to cross each other}{przecinać się}{przecinać się nawzajem.}
\dictentrytr{inverse sine}{The operation that finds which corner has a given sine value.}{arcus sinus}{operacja, która znajduje kąt o danej wartości sinusa.}
\dictentrytr{isosceles triangle}{a three-sided shape with two sides equal in length}{trójkąt równoramienny}{figura o trzech bokach, w której dwa boki mają równą długość.}
\dictentrytr{iteration}{one complete repeat of a repeating process}{iteracja}{jedno pełne powtórzenie powtarzającego się procesu.}
\dictentrytr{iterative}{using a repeated rule, where each result is used to find the next one.}{iteracyjny}{wykorzystujący powtarzaną regułę, gdzie każdy wynik jest używany do znalezienia następnego.}
\dictentrytr{iterative process}{a method that repeats a rule, using the previous answer to produce the next}{proces iteracyjny}{metoda, która powtarza regułę, używając poprzedniej odpowiedzi do otrzymania następnej.}
\dictentrytr{lattice method}{a written multiplication method that uses a grid to organise the calculation.}{metoda kratkowa}{pisemna metoda mnożenia, która używa kratek do uporządkowania obliczenia.}
\dictentrytr{leg}{one of the two shorter sides of a right-angled triangle.}{przyprostokątna}{jeden z dwóch krótszych boków trójkąta prostokątnego.}
\dictentrytr{like terms}{terms that have the same letter or letters, or no letter, so they can be added or subtracted together.}{wyrazy podobne}{wyrazy, które mają tę samą literę lub litery albo nie mają żadnej litery, więc można je dodawać lub odejmować.}
\dictentrytr{line segment}{a straight part of a line that starts at one point and ends at another.}{odcinek}{prosta część linii, która zaczyna się w jednym punkcie, a kończy w drugim.}
\dictentrytr{linear}{making a straight line when drawn as a graph}{liniowy}{tworzący linię prostą, gdy narysowany jako wykres}
\dictentrytr{lowest common multiple}{the smallest number that two or more numbers all divide into exactly}{najmniejsza wspólna wielokrotność}{najmniejsza liczba, która jest podzielna bez reszty przez każdą z dwóch lub więcej liczb}
\dictentrytr{magnitude}{the size of a number without its sign.}{wartość bezwzględna}{wielkość liczby bez jej znaku.}
\dictentrytr{major arc}{the longer of the two parts of the circumference between two points on a circle.}{łuk większy}{dłuższa z dwóch części okręgu między dwoma punktami na okręgu.}
\dictentrytr{major segment}{the larger of the two parts of a circle cut off by a chord.}{większy odcinek koła}{większa z dwóch części koła odciętych przez cięciwę.}
\dictentrytr{mapping}{matching or sending each point of one shape to a point of another.}{odwzorowanie}{przyporządkowywanie lub odwzorowywanie każdego punktu jednej figury na punkt drugiej figury.}
\dictentrytr{mass}{the amount of matter in an object, measured in grams or kilograms.}{masa}{ilość materii w obiekcie, mierzona w gramach lub kilogramach.}
\dictentrytr{maturity}{the time when a bond is paid back.}{termin wykupu}{czas, w którym obligacja jest spłacana.}
\dictentrytr{mean}{the average found by adding all the values and dividing by how many there are}{średnia}{średnia obliczana przez dodanie wszystkich wartości i podzielenie przez ich liczbę}
\dictentrytr{median}{the middle value when the data is put in order}{mediana}{środkowa wartość, gdy dane są uporządkowane}
\dictentrytr{minor arc}{the shorter of the two parts of the circumference between two points on a circle.}{łuk mniejszy}{krótsza z dwóch części okręgu między dwoma punktami na okręgu.}
\dictentrytr{minor segment}{the smaller of the two parts of a circle cut off by a chord.}{mniejszy odcinek koła}{mniejsza z dwóch części koła odciętych przez cięciwę.}
\dictentrytr{mixed number}{a number made of a whole amount and a fraction, such as 3 1/2}{liczba mieszana}{liczba składająca się z całości i ułamka, na przykład 3 1/2.}
\dictentrytr{mode}{the value that appears most often}{dominanta}{wartość, która pojawia się najczęściej}
\dictentrytr{multiple}{a number you get by multiplying another number by a whole number}{wielokrotność}{liczba, którą otrzymujesz, mnożąc inną liczbę przez liczbę całkowitą}
\dictentrytr{multiply}{to add a number to itself a given number of times}{mnożyć}{dodać liczbę do siebie określoną liczbę razy}
\dictentrytr{negative}{below zero; a number or term with a minus sign is below zero.}{ujemny}{poniżej zera; liczba lub wyraz ze znakiem minus jest poniżej zera.}
\dictentrytr{negative number}{a number less than zero, written with a minus sign.}{liczba ujemna}{liczba mniejsza od zera, zapisywana ze znakiem minus.}
\dictentrytr{non-reflex angle}{the smaller angle between two lines, measuring less than 180 degrees}{kąt wypukły}{mniejszy kąt między dwiema prostymi, mierzący mniej niż 180 stopni.}
\dictentrytr{nth term}{a rule that gives any number in a sequence from its position}{wzór na n-ty wyraz}{reguła, która daje dowolną liczbę w ciągu na podstawie jej pozycji}
\dictentrytr{number line}{a straight line with numbers marked at equal distances, used to show value and order.}{oś liczbowa}{prosta z liczbami zaznaczonymi w równych odległościach, używana do pokazywania wartości i kolejności.}
\dictentrytr{numerator}{the number above the line in a fraction, showing how many equal parts are taken}{licznik}{liczba nad kreską ułamkową, pokazująca, ile równych części zostało wziętych}
\dictentrytr{odd number}{a whole number that cannot be divided exactly by two.}{liczba nieparzysta}{liczba całkowita, której nie można podzielić dokładnie przez dwa.}
\dictentrytr{outlier}{a value far away from all the others}{wartość odstająca}{wartość, która jest daleko od wszystkich pozostałych}
\dictentrytr{parallel}{always the same distance apart and never meeting}{równoległy}{zawsze w tej samej odległości od siebie i nigdy się nie przecinające}
\dictentrytr{percentage}{a number of parts out of a hundred}{procent}{liczba części na sto}
\dictentrytr{perimeter}{the total distance around the edge of a shape}{obwód}{całkowita odległość wokół krawędzi figury}
\dictentrytr{perpendicular}{meeting at a right angle}{prostopadły}{spotykające się pod kątem prostym}
\dictentrytr{place value}{the value a digit has because of where it sits in a number}{wartość pozycyjna}{wartość, jaką cyfra ma ze względu na to, gdzie znajduje się w liczbie}
\dictentrytr{polygon}{a flat shape with straight sides}{wielokąt}{płaska figura o prostych bokach}
\dictentrytr{power}{how many times a number is multiplied by itself}{potęga}{ile razy liczba jest mnożona przez samą siebie}
\dictentrytr{prime}{a whole number greater than 1 that divides only by itself and 1}{liczba pierwsza}{liczba całkowita większa od 1, która dzieli się tylko przez siebie i 1}
\dictentrytr{probability}{how likely something is, given as a number between 0 and 1}{prawdopodobieństwo}{jak prawdopodobne jest coś, podawane jako liczba między 0 a 1}
\dictentrytr{product}{the answer when two or more numbers are multiplied}{iloczyn}{wynik mnożenia dwóch lub więcej liczb}
\dictentrytr{proportion}{a relationship where two amounts change by the same factor}{proporcjonalność}{zależność, w której dwie wielkości zmieniają się o ten sam czynnik}
\dictentrytr{quadratic}{an expression or equation whose highest power is a square}{kwadratowy}{wyrażenie lub równanie, którego najwyższa potęga jest kwadratem}
\dictentrytr{quartile}{a value marking one quarter of the way through ordered data}{kwartyl}{wartość wyznaczająca jedną czwartą drogi przez uporządkowane dane}
\dictentrytr{quotient}{the answer when one number is divided by another}{iloraz}{wynik dzielenia jednej liczby przez drugą}
\dictentrytr{radius}{the distance from the centre of a circle to its edge}{promień}{odległość od środka okręgu do jego krawędzi}
\dictentrytr{range}{the difference between the largest and smallest values}{rozstęp}{różnica między największą i najmniejszą wartością}
\dictentrytr{ratio}{a way of comparing two or more amounts, showing how much of one there is for each of the other}{stosunek}{sposób porównywania dwóch lub więcej wielkości, pokazujący, ile jednej przypada na każdą z pozostałych}
\dictentrytr{rearrange}{to change the order of a formula so a different letter is on its own}{przekształcić}{zmienić porządek wzoru tak, aby inna litera była sama}
\dictentrytr{reciprocal}{the number you multiply by to get 1, found by turning a fraction upside down}{odwrotność}{liczba, przez którą mnożysz, aby otrzymać 1, znajdowana przez odwrócenie ułamka}
\dictentrytr{recur}{to repeat the same digits forever after the decimal point}{powtarzać się}{powtarzać te same cyfry w nieskończoność po przecinku}
\dictentrytr{root}{a number that gives another number when multiplied by itself a given number of times}{pierwiastek}{liczba, która pomnożona przez samą siebie określoną liczbę razy daje inną liczbę}
\dictentrytr{round}{to replace a number with a nearby simpler one}{zaokrąglić}{zastąpić liczbę pobliską prostszą liczbą}
\dictentrytr{sample}{a smaller group chosen to stand for a larger one}{próba}{mniejsza grupa wybrana do reprezentowania większej}
\dictentrytr{scale factor}{the number every length is multiplied by when a shape is enlarged}{współczynnik skali}{liczba, przez którą mnożona jest każda długość przy powiększaniu figury}
\dictentrytr{sector}{a slice of a circle between two radii}{wycinek koła}{kawałek koła między dwoma promieniami}
\dictentrytr{segment}{the part of a circle cut off by a chord}{odcinek koła}{część koła odcięta cięciwą}
\dictentrytr{sequence}{a list of numbers that follows a rule}{ciąg}{lista liczb, która podlega regule}
\dictentrytr{similar}{the same shape but a different size}{podobny}{ten sam kształt, ale inna wielkość}
\dictentrytr{simplify}{to write something in its smallest or plainest form without changing its value}{uprościć}{zapisać coś w najmniejszej lub najprostszej postaci bez zmiany jego wartości}
\dictentrytr{substitute}{to put a number in place of a letter}{podstawić}{wstawić liczbę w miejsce litery}
\dictentrytr{sum}{the answer when numbers are added together}{suma}{wynik dodawania liczb}
\dictentrytr{surd}{a root that cannot be written exactly as a fraction, so it is left in root form}{pierwiastek niewymierny}{pierwiastek, którego nie można zapisać dokładnie jako ułamka, więc pozostaje w postaci pierwiastka}
\dictentrytr{tangent}{a line that touches a curve at exactly one point}{styczna}{prosta, która dotyka krzywej w dokładnie jednym punkcie}
\dictentrytr{term}{one part of an expression or one number in a sequence}{wyraz}{jedna część wyrażenia lub jedna liczba w ciągu}
\dictentrytr{transformation}{a change to a shape's position, size or orientation}{przekształcenie}{zmiana położenia, wielkości lub orientacji figury}
\dictentrytr{unitary}{a method that first finds the value of one part or one item}{metoda jednostkowa}{metoda, która najpierw znajduje wartość jednej części lub jednego przedmiotu}
\dictentrytr{variable}{a letter standing for a quantity that can change}{zmienna}{litera oznaczająca wielkość, która może się zmieniać}
\dictentrytr{vector}{a quantity with both size and direction}{wektor}{wielkość mająca zarówno długość, jak i kierunek}
\dictentrytr{vertex}{a corner where two or more edges meet}{wierzchołek}{róg, w którym spotykają się dwie lub więcej krawędzi}
\dictentrytr{volume}{the amount of space inside a solid shape}{objętość}{ilość miejsca wewnątrz bryły}

\end{document}
